\chapter{Architetture dei Modelli e Setup Sperimentale}
\label{ch:architetture}

\section{Implementazione della PCA come baseline}
\label{sec:setup_pca}

Al fine di stabilire un \emph{benchmark} quantitativo rigoroso per valutare le performance delle reti neurali, la pipeline sperimentale prevede l'implementazione iniziale della \emph{Principal Component Analysis} (PCA). Questo modello funge da base comparativa per determinare il limite massimo di compressione lineare raggiungibile sui dati.

La PCA è stata addestrata sul \emph{Training Set} linearizzato, composto dai vettori $\mathbf{x} \in \mathbb{R}^{65341}$ rappresentanti le componenti triangolari inferiori delle matrici di correlazione. Fissata una dimensione target $K$ per lo spazio compresso, l'algoritmo estrae gli autovettori corrispondenti ai $K$ autovalori dominanti della matrice di covarianza globale del set di addestramento. Le performance di ricostruzione, misurate tramite l'Errore Quadratico Medio (MSE) sulle matrici del \emph{Test Set}, costituiscono la soglia che i modelli di \emph{deep learning} devono teoricamente eguagliare (nel caso lineare) o superare (nel caso non lineare).

\section{Design e iperparametri dell'Autoencoder Lineare}
\label{sec:setup_linear_ae}

L'Autoencoder Lineare (Linear AE) è stato progettato per fungere da ponte analitico tra l'apprendimento profondo e l'algebra lineare. Per garantire la perfetta replicabilità del sottospazio della PCA, l'architettura è stata sottoposta a vincoli ingegneristici stringenti:
\begin{itemize}
    \item \textbf{Struttura:} Un singolo \emph{layer} denso per l'Encoder ($\mathbb{R}^{65341} \to \mathbb{R}^K$) e uno speculare per il Decoder ($\mathbb{R}^K \to \mathbb{R}^{65341}$).
    \item \textbf{Funzioni di attivazione:} Rigorosamente identità (lineari) in ogni nodo della rete.
    \item \textbf{Assenza di Bias e Regolarizzazione:} I termini di \emph{bias} sono stati disabilitati in tutti i \emph{layer} ($\mathbf{b} = 0$) e non è stata applicata alcuna penalizzazione di tipo \emph{weight decay} (L2) per non alterare la magnitudo della scomposizione spettrale appresa.
\end{itemize}

L'ottimizzazione dei pesi è stata condotta minimizzando esclusivamente la funzione di costo spaziale MSE. Come prassi di \emph{sanity check}, gli iperparametri di addestramento (in particolare la scelta dell'ottimizzatore, come Adam, e il \emph{Learning Rate}) sono stati calibrati in modo iterativo per garantire che la \emph{loss} del Linear AE convergesse asintoticamente sull'orizzonte esatto dettato dalla PCA sul set di addestramento, validando l'assenza di minimi locali sub-ottimali.

\section{Design e iperparametri dell'Autoencoder Non-Lineare}
\label{sec:setup_deep_ae}

Per superare la rigidità geometrica dei modelli precedenti e catturare le interdipendenze gerarchiche e asimmetriche del mercato reale, è stato sviluppato un Autoencoder Profondo (AE). L'inclusione di strati multipli e funzioni non lineari dota il modello della flessibilità necessaria per piegare lo spazio latente e fittare la \emph{manifold} complessa dei dati finanziari.

L'architettura è strutturata con un \emph{design} a imbuto (simmetrico):
\begin{itemize}
    \item \textbf{Encoder e Regolarizzazione:} L'infrastruttura inizia con il \emph{layer} di ingresso ($\mathbb{R}^{65341}$), seguito da strati nascosti (\emph{hidden layers}) a dimensionalità decrescente, fino al livello di \emph{bottleneck} ($\mathbb{R}^K$). Per forzare la rete ad apprendere rappresentazioni più robuste e distribuite delle dinamiche di mercato, mitigando contestualmente il rischio di memorizzazione mnemonica (\emph{co-adaptation}), è stato interposto un \emph{layer} di \emph{Dropout} tra gli strati dell'Encoder.
    \item \textbf{Decoder:} Struttura speculare che riceve il vettore latente $K$-dimensionale e lo proietta nuovamente alla dimensione originale attraverso un'espansione simmetrica.
    \item \textbf{Funzioni di attivazione:} Per tutti gli strati nascosti è stata adottata la funzione di attivazione LeakyReLU (con parametro di pendenza $\alpha$ fissato a 0.01) per scongiurare il problema del \emph{vanishing gradient} e garantire un flusso di gradiente non nullo anche per valori di input negativi. Per l'ultimo \emph{layer} di output del Decoder è stata sistematicamente impiegata l'attivazione Tanh (tangente iperbolica), una scelta architetturale atta a mappare rigorosamente i valori ricostruiti nel dominio continuo $[-1, 1]$, coerente con la natura geometrica dei coefficienti di correlazione di Pearson.
\end{itemize}

Al fine di prevenire l'\emph{overfitting} e selezionare la configurazione ottimale, il protocollo di addestramento prevede l'esecuzione per un numero prefissato di epoche, monitorando costantemente il \emph{Validation Set} (generato tramite la tecnica del \emph{random shuffle}). La strategia di salvataggio dei pesi (\emph{Model Checkpointing}) è impostata per conservare esclusivamente il modello in corrispondenza del valore minimo assoluto della \emph{Validation Loss}. 

Grazie all'omogeneità statistica garantita dal campionamento casuale, che annulla il problema del \emph{distribution shift} temporale, il Deep AE dimostra una capacità espressiva eccezionale: la \emph{Training Loss} e la \emph{Validation Loss} risultano stabilmente decrescenti, scendendo al di sotto del limite rigido fissato dal \emph{benchmark} della PCA per medesimi valori di $K$. Questo successo empirico certifica la capacità della rete di superare i limiti dell'algebra lineare, generalizzando efficacemente sulle dinamiche non lineari del mercato.

Tuttavia, l'eccellente potere ricostruttivo del Deep AE porta alla luce un limite strutturale di diversa natura: l'incapacità generativa. Trattandosi di un modello puramente deterministico, l'Encoder mappa l'input in vettori latenti $\mathbf{z}$ puntuali, creando uno spazio disaggregato e privo di organizzazione probabilistica. Estrapolare nuovi scenari campionando vettori casuali da questo dominio produrrebbe matrici del tutto prive di validità macroeconomica. Tale ostacolo sancisce definitivamente la necessità di transitare verso il paradigma variazionale (VAE), progettato esattamente per superare il determinismo e abilitare la simulazione stocastica.

\section{Design e iperparametri del VAE}
\label{sec:setup_vae}

Il Variational Autoencoder (VAE) rappresenta l'apice ingegneristico della pipeline, deputato non solo alla compressione, ma alla generazione stocastica. L'architettura espande la topologia del Deep AE trasformando lo strato di \emph{bottleneck} deterministico in un modulo di inferenza probabilistica.

\begin{itemize}
    \item \textbf{Rete di Inferenza (Encoder):} Al termine degli strati nascosti, la rete si biforca in due layer paralleli che emettono, rispettivamente, il vettore delle medie $\boldsymbol{\mu}$ e il vettore delle varianze logaritmiche $\log(\boldsymbol{\sigma}^2)$ della distribuzione latente, per garantire la stabilità numerica.
    \item \textbf{Sampling e Reparameterization:} Lo strato latente non riceve un input fisso, ma un campione estratto dinamicamente tramite il \emph{reparameterization trick} $\mathbf{z} = \boldsymbol{\mu} + \boldsymbol{\sigma} \odot \boldsymbol{\epsilon}$, dove $\boldsymbol{\epsilon} \sim \mathcal{N}(\mathbf{0}, \mathbf{I})$.
    \item \textbf{Rete Generativa (Decoder):} Il Decoder speculare riceve $\mathbf{z}$ e procede alla ricostruzione.
\end{itemize}

Il bilanciamento tra l'errore di ricostruzione (MSE) e il vincolo topologico imposto dalla Divergenza di Kullback-Leibler (KL) è l'iperparametro critico del sistema. L'addestramento modella uno spazio continuo che permette, in fase di test, di estrarre nuove matrici sintetiche campionando puramente dalla normale standard $p(\mathbf{z})$ senza necessitare di input originari.

\section{Metriche di valutazione e criteri di convergenza}
\label{sec:metriche_convergenza}

L'addestramento dei modelli neurali è stato governato monitorando un set di metriche oggettive e applicando precise policy di ottimizzazione per massimizzare la stabilità dell'apprendimento e la capacità di estrapolazione fuori campionario:

\begin{enumerate}
    \item \textbf{Metriche Spaziali Rigide:} La valutazione quantitativa puntuale si basa sul \emph{Mean Squared Error} (MSE), impiegato direttamente come funzione di costo e indicatore primario. A questo si affiancano il \emph{Mean Absolute Error} (MAE), più resistente agli \emph{outlier}, e la Norma di Frobenius, che quantifica l'energia complessiva residua tra la matrice empirica originale e la sua approssimazione ricostruita o generata.
    
    \item \textbf{Ottimizzatore e Learning Rate Policy:} È stato impiegato l'algoritmo di ottimizzazione Adam, rinomato per la gestione dinamica del momento del gradiente. Per facilitare il superamento del limite lineare sul \emph{Training Set} e agevolare la discesa verso il minimo globale, è stato implementato uno \emph{scheduler} per il \emph{Learning Rate} basato sulla tecnica del \emph{Cosine Annealing}. Questa strategia riduce il tasso di apprendimento seguendo il decadimento morbido e continuo della funzione coseno, permettendo esplorazioni ampie nelle fasi iniziali e un'ottimizzazione fine, priva dei bruschi salti tipici degli approcci \emph{step-based}, man mano che la rete si avvicina al \emph{plateau} di convergenza.
    
    \item \textbf{Monitoraggio della Convergenza:} La gestione a epoche (\emph{epochs}) tramite \emph{mini-batches} è stata affiancata dall'impiego del \emph{Validation Set} come arbitro indipendente per il \emph{Model Checkpointing}. L'algoritmo valuta l'intero ciclo di epoche prefissato e salva i pesi architetturali esclusivamente in corrispondenza del valore minimo assoluto raggiunto dalla \emph{Validation Loss}. Questo protocollo garantisce che le prestazioni valutate in sede finale sul \emph{Test Set} costituiscano la stima più imparziale e robusta possibile per gli scenari futuri di \emph{asset allocation}.
\end{enumerate}